\quellseite{430}
\chapter[§ 31. THE DISTRIBUTION OF THE MATERIAL]{§ 31. THE DISTRIBUTION OF THE MATERIAL OVER THE DIFFERENT EDUCATIONAL AGENCIES AND THEIR SEVERAL GRADES}

Compare the literature given in §§ 20 and 30. *F. Zange, pp. 64–142. A. Eckert, pp. 62–78. J. Berndt, pp. 68–73. R. Kabisch, pp. 157–181. R. Seyfert, \emph{Versuch eines Lehrplans fuer den Religionsunterricht}, 1890. J. H. Schueren, \emph{Gedanken ueber den Religionsunterricht}, 1900. W. Armstroff, \emph{Die einheitliche Gestaltung des Religionsunterrichts in Kirche und Schule}, 1896. H. Brammer, \emph{Neue Bahnen fuer den Religionsunterricht}, 1900–1903. *A. Reukauf, \emph{Zur Lehrplantheorie der geschichtlichen Stoffe im Religionsunterricht der Volksschule}, 1906. Baumann, \emph{Lehrplan f. d. ev.-luth. Religionsunterricht der achtklassigen Volksschule}, 1906. S. Bang, \emph{Lehrplan f. d. 2–8 klassige einfache Volksschule}, 1906. A. Reukauf, \emph{Didaktik des evang. Religionsunterrichts}, 1906. *G. W. Pease, \emph{An Outline of a Bible School Curriculum}, 1909.

So far as the home is concerned, a distribution of the catechetical material is out of the question. While the home, as the child grows up, becomes largely the custodian of blessings imparted to the child in successive stages, any communication of religious truth in the childhood stage is of necessity spontaneous and occasional. Of what radical and decisive moment the infant period is for the later development and the formation of character, has already been shown on pages 245 ff., where also the needed hints as to the material to be employed are given. Here we merely desire to emphasize the point that the education dealing with the infant dare not disturb the quiet sanctuary of the awakening inner life; its aim will be above all to clarify, to guide, to achieve results through the living example. Here above everything the pedantic tone of the schoolmaster should be avoided and all cold reflection likewise! Cautious and chaste nurture of the nascent life imbedded within through Baptism should be the aim.

When we come to speak of the second educational agency, however, the school, a distribution of the material is necessary.

\quellseite{431}
We concern ourselves first of all with the regular parish school. Since Biblical History is to be taught in all its grades, the principle governing the distribution of this material should be clearly understood first of all. For a long time, the curriculum was devised on the \emph{concentric circle plan}, according to which originally all Bible stories were gone through with all the classes year after year (the Old Testament in the summer, the New Testament in the winter). Later this was changed in that a few, easily comprehended stories were taken up in the Lower Course, a review of these, with a few added ones, was undertaken in the Intermediate Course, and the full number of stories contained in the Biblical History was reached in the Upper Course. This method of distribution was made an object of attack by Ziller and his followers who advocated with all their might a curriculum based on the successive culture epochs, \emph{kulturgeschichtliche Stufen}. According to their plan the children should be made to live over again in their own experience the history of human development from the Patriarchs to the life of faith typified by the Evangelical Church. In pursuance of this plan, a preliminary course, ranging over two or three years, was arranged, with tales and Robinson Crusoe or with religious subjects as material. This was followed by the history of the Patriarchs in the third school year, the story of Moses and the Judges in the fourth, the story of David and Solomon in the fifth, the story of Jesus together with selections from the Prophets in the sixth, the Acts of the Apostles together with parts of the Epistles in the seventh, and the history of the Reformation together with a final resume of the Catechism in the eighth. Against the curriculum patterned upon concentric circles, as originally devised, the argument was advanced that, in carrying it out, too much was gone over in one year, which prevented thoroughness and for this reason made impossible the inner experience and aroused no speculative interests. Against its first form and its subsequent
\quellseite{432}
modifications the argument was raised that, with a repeated treatment of the same stories, ennui resulted for the children—an outcome unfavorable to the arousing of the religious and sympathetic interests. These arguments really disclose the weaknesses of the curriculum based upon the concentric circle plan. But the curriculum based on the culture epochs is not qualified to take its place. A crushing argument against it is the untenableness of the underlying idea, which, ignoring the fact of Baptism, does not introduce the child to religious materials until the third or fourth school year. It can also be urged that, without a repeated treatment of the subject matter of Biblical History, it is impossible to bring about such conversance with it as the young stand in need of upon entering life, even the socalled immanent review\footnote{By immanent review we mean the review connected with the treatment of the new material; I employ immanent review either when I proceed from the known material to discover unknown truths, or when I combine the new material with the old and fuse them into a unit. Compare the subject of “penetration”, § 32.} not being sufficient for the purpose. What is, however, eminently worthy of attention is the effort to present the Bible stories in large groups, for only in this way a vital penetration into the history of the several biblical characters, is made possible and, likewise, the arousing of a many-sided interest, sympathetic, religious, moral. Accordingly, a curriculum that possesses the advantages of the two mentioned, without being marred by their faults, will be found to be the most efficient one. For this reason simple, but stirring stories chosen from the New Testament as well as the Old, should be introduced already in the Lower Course; with the proviso that the stories selected cover not a great number of periods but only one or two, and, in any case, constitute a natural unit. In the Middle Course there should be a review of the material previously gone through; but these stories, now known, should at this time be placed under a higher aspect and not be
\quellseite{433}
treated at the same length, while the new stories should by far preponderate. But whatever material is taken up for treatment should constitute a compact group or a number of compact groups. All the lesson material hitherto used will be coordinated in the Upper Course into a coherent Sacred History. In the Lower Course the selection of material should conform as much as possible to the church year; in the Middle Course at least the festive stories will be reviewed in due time, while the church year receives no further consideration in the Upper Course than is necessary to read and explain the pericopes.

Carrying out the principles here laid down, we offer the following curriculum, sketched upon the supposition that the three grades of the lower course have a common instructor, likewise the middle course of two grades and the upper course, also consisting of two grades. The material assigned to the first or second grade of each course is so shifted that the pupils entering in 1927 and those who have entered in 1926 begin with the material selected for the second grade and then, in the second year take up the material designed for the first grade. Since, in the lower course, all instruction should be fused into an organic unit with Biblical History, we show at the hand of several examples from the lower course how such fusion is to be accomplished. At the same time we refer to our \emph{Wartburg Lesson Helps} and our \emph{Biblical History} (Chicago, 1918) where this fusion has been effected.

\emph{First or Primary Course} (comprising three school years, age 6–9). \emph{First Grade:} (1) How God created the world. A (=first year): Ps. 115:3; B (=second year): Ps. 104:24; C (=third year): Ps. 33:9; Morning prayer; Catechism: First Article without the explanation.—(2) How God made man happy. A: Ps. 118:1; B: 1 John 4:19; Catechism: First Commandment without explanation; Hymn: Awake, my heart rejoicing, 1st stanza; C: 1 John 5:3; Catechism: First Commandment with explanation; Hymn: Praise to the Lord, the Almighty, stanzas 1 and 2.—(3) How man sinned. A: Ps. 5:4; Prayer: Abide in grace, Lord Jesus, 1st stanza; B: Prov. 1:10; Prayer: Lord Jesus, who dost love me; C: Ps. 37:37; Catechism: Conclusion of the Commandments without explanation; Hymn: Abide in grace, Lord Jesus, stanzas 1, 5, 6.—(4) The Saviour
\quellseite{434}
came into the world. A: Luke 2:10–11; Hymn: As each happy Christmas; B. John 3:16; Hymn: Praise God the Lord, ye sons of men, stanzas 1, 2, 8. C: Luke 19:10; Hymn: Good news from heav'n, stanzas 1–4.—(5) Wise men from the East came to the child Jesus. A: Matt. 7:7; Hymn: Brightest and best of the sons of the morning; B: Prov. 23:26; Hymn: Thou, whose almighty word; C: 1 Tim. 2:4; Acts 4:12; Hymn: Thou whose almighty word; Catechism: Go ye and teach all, etc.—(6) How the child Jesus was saved from great danger. A: Ps. 34:7; Prayer: Lord Jesus, who dost love me; B: Ps. 91:11; Hymn: If thou but suffer God to guide thee, 1st stanza; C: Ps. 23:3; Hymn: Thy way and all thy sorrows, stanzas 1 and 2.—(7) The child Jesus is brought to the temple. A: Luke 11:28; Catechism: Third and Fourth Commandments without explanation; B: Ps. 26:8; Hymn: Blessed Jesus, at thy word, 1st stanza; C: Third and Fourth Commandments with explanation.—(8) How Jesus blessed Peter’s draught of fishes. A: All depends on our possessing God’s true love, and grace, and blessing; B: Ps. 127:1; Hymn: Take thou my hand, O Father, 1st stanza; C: Matt. 10:37; Hymn: Jesus, still lead on, 1st stanza; Abide in grace, 4th stanza.—(9) How Jesus saved His disciples during a storm. A: Ps. 50:15; Hymn: Ever is a peril near me, 1st stanza; B: Matt. 28:18; Hymn: Abide in grace, 6th stanza; C: Rom. 8:31; Hymn: How the wind in fury blind, 1st stanza.—(10) Jesus feeds the hungry multitude. A: Come, Lord Jesus, be our guest, etc. B: Ps. 103:13; Catechism: Prayer before meat; C: Ps. 145:15. 16; Hymn: If thou but suffer, stanzas 1–3; Catechism: Prayer after meat.—(11) Jesus heals the impotent man at the pool of Bethesda. A: Luke 1:37; B: Ps. 118:8; C: Ps. 103:8; Hymn: Thy way and all thy sorrows, stanza 6.—(12) How Jesus comforted a sorrowing father by restoring his daughter to life. A: I am the resurrection and the life; B: Mark 5:36 (Fear not; only believe); C: Ps. 25:3; Hymn: When children young and tender, 1st stanza.—(13) How Jesus restored a son to his mother. A: Rom. 12:5; B: 2 Tim. 1:10; C: John 5:28. 29.—(14) Jesus rejoices with the rejoicing and helps them out of their difficulty. A: Rom. 12:5; B: Ps. 34:8; Hymn: I am Jesus’ little lamb, 1st stanza; C: Ps. 77:14; Hymn: Thy way and all, 4th stanza.—(15) Jesus desires that little children should be brought to Him. A: Mark 10:14; Hymn: The truest Friend abides in heaven; B: Hymn: Jesus now thine own forever; C: Hymn: I am Jesus’ little lamb, stanzas 1–3.—(16) How Jesus entered Jerusalem
\quellseite{435}
on Palm Sunday. — (17) How Jesus prayed before His suffering. — (18) How Jesus was betrayed by one of His disciples. — (19) How Jesus died for us on the cross. — (20) How Jesus was buried on Good Friday. — (21) How the risen Saviour on Easter day appeared unto His disciples (John 20). — (22) How Jesus ascends into heaven in order to be with us always.

\emph{Second Grade} (1) Why Abraham, the man of faith, left his home. (2) How Abraham preferred to be separated from Lot rather than live in strife and contention. (3) How Abraham proved to be a friend in need. (4) What pleased God most in Abraham. (5) Abraham invites three guests and thus receives a new promise. (6) Abraham intercedes for the people of Sodom. (7) How Sodom and Gomorrah were destroyed because of their sin. (8) God gives Abraham a son, but Abraham loves God more than his son. (9) God provides a wife for Isaac. (10) Abraham dies and is buried by his son. (11) Jacob deceived his father and his brother. (12) Jacob fled from his brother. (13) Jacob saw the heavens open in his dream. (14) Of Jacob’s life in the home of Laban. (15) Jacob flees from Laban. (16) How Jacob meets his brother Esau. (17) Of Joseph, Jacob’s favorite son. (18) Of Joseph’s life in the house of Potiphar. (19) What happened to Joseph in prison. (20) How Joseph was exalted and made lord of all Egypt. (21) Joseph’s brothers come to Egypt. (22) Joseph’s brothers come to Egypt for the second time. (23) Joseph makes himself known to his brothers and treats them kindly. (24) Jacob again sees his son Joseph. (25) Jacob blesses his children and dies. (26) Joseph’s death. — To these, the festival stories are again added at the proper time, so that at least the most important facts of the life of Christ are touched upon also in this grade.

\emph{Third Grade:} In the third year the stories given above under grade 1 are again gone through. As a rule, one Bible story is gone through for each week of instruction; sometimes two weeks will be necessary. Allowing for the proper number of weeks for the review, which must set in after groups of stories belonging together have been taught, the usual number of weeks of school (about 36) is just sufficient for the material given above. The Bible verses, prayers, and hymns, which must be brought in relation to the Bible story should be assigned for study also on other days of the week, but never before light has been cast upon them from the Bible story to which they belong, and their recitation must always be introduced by questions bringing
\quellseite{436}
out their relation to the Bible story. The most important memory material for the three grades of the first course is given in the following, and should be brought in connection with the instruction in Bible History in the manner suggested above (compare \emph{Wartburg Lesson Helps, Primary Grade}):

\emph{First Year or Grade:} Bible verses: Ps. 115:3; 1 John 4:19; Ps. 5:4; 133:1; Matt. 5:9; Ps. 37:5; 37:37; 50:15; John 3:16; Prov. 23:26; Luke 11:28; Mark. 10:14; Prov. 11:1; John 1:29. Prayers: Morning prayer, prayer before meat, prayer after meat, evening prayer. Hymns: As each happy Christmas; Good news from heav’n, stanzas 1, 2; O sacred head now wounded, 6th stanza; Abide in grace, Lord Jesus, 1st stanza.—For the second year or grade: Bible verses: 1 John 5:3; Ps. 26:8; 1 John 4:16; Gen. 17:1; Ps. 33:9; 36:5; 34:7; 103:2; Luke 19:10; 1 John 1:7; Ps. 51:10; Gen. 39:9; Prov. 1:8; Matt. 5:8; Eph. 4:25; 1 Cor. 6:14; Romans 14:8; 8:28; Ps. 118:8; Ps. 91:11; 127:1; Mark 5:36; Romans 12:5. Prayers and Hymns: Now to gain a night’s repose, stanzas 1, 2; Ah, dearest Jesus, holy child, make Thee a bed soft, undefiled, within my heart, that it may be a quiet chamber kept for Thee; Awake my heart rejoicing, 1st stanza; Praise God the Lord, ye sons of men, stanzas 1, 2; Thou whose almighty word; If thou but suffer God to guide thee, stanzas 1, 3; I could not do without thee, 2nd stanza; I am Jesus’ little lamb, stanzas 1, 2; Abide in grace, Lord Jesus, stanzas 5, 6; Jesus, still lead on, 1st stanza; Glory be to Jesus, who in bitter pains, stanzas 1, 2; My faith looks up to Thee, stanzas 1, 2. Catechism: The Ten Commandments without explanation.—For the third year or grade: Bible verses: Ps. 119:105; John 5:39; Ps. 23:1; Ps. 23:2; Ps. 23:4; Ps. 104:24; Ps. 103:8; Gen. 8:22; Gen. 32:10; 8:21; John 14:6; Matt. 28:20; 1 Tim. 2:4; 2 Tim. 1:10; 2 Peter 3:9; 1 Peter 5:9; 1 Tim. 5:4; Hebrews 13:7; 1 Peter 3:9; Isa. 58:7; 1 John 2:15; Lev. 19:2; Gal. 6:7; Ps. 68:20; 2 Cor. 5:10; Isa. 28:29; John 5:28, 29. Prayers and Hymns: The Lord’s Prayer, Good news from heav’n, stanzas 1–6; Praise God, from whom all blessings flow; Jesus, Thy blood and righteousness; Jesus, now Thine own for ever, 1st stanza; Blessed Jesus, at Thy word, 1st stanza; Now the day is over; Glory be to Jesus; I am Jesus’ little lamb; He lives, my Lord has left the grave; O Holy Spirit, enter in, 1st stanza; Catechism: The first three Chief Parts without explanation. The explanation may be added to a few of the commandments. In each grade, the memory material previously learned is to be reviewed, and again brought into connection with the Bible story to which it belongs.

\quellseite{437}
\emph{Second or Intermediate Course} (comprising 2 school years or grades, age 10 and 11). \emph{First year or grade:} Bible History: History of creation and early history of man; Review of the history of the Patriarchs; stories of Moses; stories of Joshua; stories of the Judges; stories of Saul, David, and Solomon (the difficult stories may be omitted in this grade). Where time permits, a few stories of Elijah may be added. Review of the festival stories.—Hymns: Praise God the Lord, ye sons of men; Praise ye the Lord in simple joyous measure; My faith looks up to Thee; All hail the pow’r of Jesus’ name; Rest of the weary; Jesus loves me; Jesus lives, thy terrors now; If thou but suffer God to guide thee; Thy way and all thy sorrows; Now thank we all our God; Praise to the Lord, the almighty. Bible verses: One half of the verses marked with two stars in our “Catechism with Explanation”, and review of those learned in the First or Primary Course.—Catechism: The First and Second Chief Parts with explanation.

\emph{Second Year or Grade:} Biblical History: Stories from the life of Christ in groups f. i.: John the forerunner of Christ; The Savior’s youth; Jesus’ public ministry; Jesus the helper of the sick; Jesus the Savior in sorrow and death; Jesus the Savior of sinners; Jesus the teacher of the people; Christ’s success; Jesus disputes with the Pharisees; Jesus’ suffering and victory. Hymns: O how shall I receive Thee, stanzas 1, 2; My dear Jesus, I’ll not leave; Come, follow me, the Savior spake; My life is hid in Jesus, stanzas 1, 2; Jesus sinners doth receive; O Thou that hear’st when sinners cry; O sacred head now wounded, stanzas 1, 8, 10; Abide in grace, Lord Jesus.—Bible verses: One half of the verses marked with two stars in our catechism.—Catechism: The first three Chief Parts with explanation. Morning, evening and table prayers.

N. B. As in Course 1, the material given above is to be brought into connection with the instruction in Bible History; therefore the catechist must plan beforehand how this can best be done, in order to enable him to assign the memory material in the proper order. Compare \emph{Wartburg Lesson Helps, Second Course}. Intermediate Grade, and the Bible History published by the author.

\emph{Third or Upper Course} (comprising 2 years or grades, age 12 and 13). \emph{First Year or Grade:} Biblical History: The history of salvation in the Old Testament to Solomon and the division of the kingdom. History of salvation from the division of the kingdom to the restoration (compare also the Appendix to the Old Testament portion of the Bible History); with special emphasis on the labors and books
\quellseite{438}
of the Prophets.—Bible Reading.—Bible Study of the Old Testament.—Catechism: In this grade special periods for Catechism are set aside in which the Ten Commandments and the First Article are gone through according to the Catechism with explanation. Hymns and Bible verses are brought into connection with such instruction as much as possible, also Bible reading.—Bible verses: The verses marked with a cross in our catechism.—Hymns: O sacred head, now wounded; O Holy Spirit, enter in; Jesus Christ, my sure defense; A mighty fortress is our God.

\emph{Second Year or Grade:} Biblical History: The history of salvation in the New Testament to the ascension of Christ; The outpouring of the Holy Ghost and the establishing of the Christian Church by the Apostles, with special emphasis upon the Epistles. Sketches from Church History (see the appendix to the Catechism with explanation). Bible Reading. Introduction to the New Testament.—Catechism: Explanation of the Second and Third Article, and of the Third Chief Part. It may, however, be desirable, that the first three Chief Parts are completely treated in the second as well as in the first year. Memorizing of the Fourth and Fifth Chief Parts, also of the Part on Confession and Absolution, if this has not been done in the first year. Bible verses: Review of all Bible verses, and memorizing of longer portions of Holy Scriptures. Compare Ch. 27.—Hymns: Why, my soul, thus trembling ever; Who knows how near my life's expended; With all my heart I love Thee, Lord; What our Father does is well; The Church's one foundation.

Such distribution of material being calculated for a regular parish school, the question arises how the material is to be distributed where the establishment of a regular parish school has proved inexpedient. If there is a partial parish school, so arranged that children from the sixth to the ninth year are in attendance, that share of the material which, according to the preceding sketch, is allotted to the First Course of a regular parish school, falls to such a partial school, special attention to be paid to the simple, wholesome juvenile hymns, as found in the author's For Beginners (translated by H. Brueckner), or in a good juvenile hymnal. The material for the Middle Course is thereupon disposed of in the combination
\quellseite{439}
school (Summer, Sunday, Saturday), while the most important parts of the material assigned to the Upper Course of a regular parish school are gone over during catechumenal instruction, especially a survey of Biblical History, a number of Bible passages with explanation, and the essential facts relative to the church year, the hymn, and the order of service. If there is a partial parish school so arranged that children from the twelfth to the thirteenth year attend it, the combination school should assume in the preceding years the most important elements of the material assigned to the Lower and Middle Courses while that assigned to the Upper Course now falls to the one-course parish school.

If success has crowned the effort to win the public school authorities for the plan to allot some of the time claimed by the secular school to a representative of the Church trained in pedagogics, the material intended for the three courses of the regular parish school will constitute that for the instruction in religion thus made possible, while the material for catechumenal instruction will be the same in that case as that allotted to catechetical instruction in connection with a regular parish school.

Should, however, anyone draw the conclusion that the adjustment above proposed obviates the necessity for a regular parish school, we reply again: Without a regular, well conducted parish school, attended by all the children of the parish, many things are taught with too much haste and, therefore, cannot be made as productive for the whole man as they should; without it, especially in view of the frequent absence of the Christian spirit from the home, religious habit—a most important factor—does not come into its own; without it, the secular branches are not taught in Christ’s spirit, with the result that two conflicting views of the world compete for mastery in the soul, or even that the one prevailing in the public school may be victorious (cf. § 20). Notwithstanding,
\quellseite{440}
we concede that this perplexing feature can be, partly at least, removed in the same degree in which the Church succeeds in penetrating the home atmosphere and in securing for the catechetical instruction given during public school hours not only scientifically trained men, but also Christian characters of the Lutheran type. These latter, however, “do not fall from heaven”, but are the result of training, true as it is that they are a gift of God. Only the best human material produced by the Church is good enough for a function so fraught with possibilities for the future as this. If the ideal, a regular parish school, lies beyond the horizon of the attainable, let there be speedy, energetic, and joint efforts to make the best of the substitutes proposed by us. No complaints, no paper resolutions will here avail, but only purposeful, unwearying labor and prayer for an effective eye salve for oneself, the congregation, and the leading men of the Church. If despite all efforts and prayers, the obstacles in the way of a parish school, complete or partial, cannot be surmounted, let that be effected which is practicable in all circumstances, even in the mission field,—the combination of Summer, Sunday, and Saturday Schools with additional hours of instruction for catechumens; and let this opportunity, in a purposeful manner, be turned to account!

\trenner

\emph{The Sunday School}, in accordance with the mental development of the pupils, is divided into three departments of which the first comprises the pupils of infant age, the second those in the childhood stage, the third the youth of both sexes. Conditions render it advisable to assign all the children up to the seventh year to the First, or Beginners’, Department. The Second Department, comprising the children from the eighth to the thirteenth year, is divided into three, or at least two, grades. And the confirmed constitute the Third, or Senior, Department. Where there is no regular parish school
\quellseite{441}
or one not patronized by all the children of the parish, the Second Course requires the most thorough attention. That Biblical History, as also in the first, constitutes the principal material for instruction, is self-evident. “In all courses, the Bible story is fundamental; for the Bible story reveals the acts of God, in which our salvation, divinely wrought, is comprehended” (Schueren, Gedanken etc., p. 9; cf. above pp. 250, 252 f., 260). Scripture passages, the Catechism, and the hymn should be connected with it for the purpose of summarizing whatever truth the Bible story contains; but first rank must be assigned to Biblical History. Separate instruction in the Catechism, moreover, would be out of the question in view of the teaching force at the command of the congregation under present circumstances.

\emph{The Second Course—Childhood}. That a continuous instruction in Biblical History in which the narratives of the Bible are taken up seriatim from the creation of the world to the missionary journeys of Paul, is not the province of the Sunday School, follows from the fact that its pupils belong to the most divergent stages of life and development. Some selection has to be made. That could be done in such a way that the material designed for the regular parish school, with a few modifications, is allotted to the Second Course of the Sunday School, especially since the pupils of this department, as of the parish school, are about of the same age and require to be divided into Primary, Intermediate, and Junior grades. That done, each grade would have to finish a certain task calculated for two years, with the result that the Second Course would be divided into three independent units. In this way the educational facilities would best adjust themselves to the mental development. These reasons have prompted Dr. Schmauk, in his \emph{Graded System of Sunday School} to assign to the children from ten to fourteen years Bible Story, Bible Readings, Biblical History, Bible Geography, and Bible
\quellseite{442}
Biography as appropriate material, which has largely been handled in a masterly manner. Where for each one of these grades there is a separate teacher who knows his business and never forgets that he has fulfilled his task not merely by introducing the lesson material into the concept world of his pupils, but has, at the same time, stirred their emotional and volitional life, we believe the adoption of the lesson plan of the parish school (with modification) or that proposed by Schmauk, or of some similar one, quite practicable and promising of good results. But the question remains whether those postulates in most cases really exist. Not seldom the Sunday School uses the same room as the congregation; in that case all its members, at all events, those of the Second Course, which we now have in mind, have to be instructed in a room not seldom without partitions, often very crowded. It meets but once a week for the space of one hour, in contrast to a parish school whose three grades, even when crowded under one teacher into one room, yet receive two hours’ instruction a week in Biblical History. Above all, while there are a number of assistants of both sexes in the Sunday School, these are frequently rather helpless while there is usually but one competent, purposeful leader, unable in the short time at his disposal, to teach all three grades himself. To be sure, the teaching force can and should be improved (while the summer sessions of the teachers’ institutes, some of which are Lutheran, are not likely to be attended by most of the teachers, there remains the weekly teachers’ meeting by way of preparation for Sunday); notwithstanding, the postulates of those lesson plans, which are necessary to success, will be realized but seldom. How can a teaching force such as we have at our command, independently present the several Bible stories, adapt them to the different grades, elaborate from the individual story the inherent truth for faith and life, and train the class for mental co-operation? In these
\quellseite{443}
circumstances instruction in Sunday School is likely to be a mere cramming of the external events of the story; but there will be no arousing of the sympathies and of religious and moral promptings—so important for development of the inner life (cf. pp. 222, 224), and through it, a stirring of the will.

This was felt by Hunt; for this reason he deems it important that the pastor should address the whole school every Sunday, “to impress a spiritual truth upon the hearts of the pupils” (in \emph{The Lutheran Sunday School Handbook}, p. 62). True, only the remedy is not adequate to the evil. The school having been occupied with divergent, disconnected material, the pastor cannot link such “spiritual truth” to the subject matter with which the several classes had been occupied; there is no effective turning to account of the latter; the “spiritual truth” presented by the pastor is felt by the class to be merely superimposed from without, and the intended effect cannot be realized. However, the needed stimulation of the intellectual, sympathetic, and above all, of the ethical and religious interests, and the consequent prompting of the will are best effected where the whole school, or at least the Second Course (age 8–12) considers the same story on one and the same Sunday, and where this has been related and presented by the leader first, so presented—according to the intuitive principle and with reference to the Catechism, Scripture, and hymn material—that it is not only understood by the whole department, but that the religious truths imbedded in the story are also developed from it for direct application. But does the gradual development, so important for mental progress and certainly entitled to consideration, in this way really come into its own? This is the case indeed; for (1) in contrast to the former lesson plans of the International Sunday School Association, it is only for the Second Course that one and the same Bible story is required for the same Sunday; (2) an efficient leader will so formulate his presentation
\quellseite{444}
that the members of the Primary Grade can understand it and that at the same time it will not become dull and insipid for the pupils of the Junior Grade; and (3) the three grades of the Second Course are given a printed lesson help, in which the story under consideration with the related material from Scripture, the Catechism and the Church Book is analyzed for each class separately, and careful attention is paid to the progressive stages of the children’s mental development (cf. Wartburg Lesson Helps, Second Course).

Where all three grades of the Second Course always have the same story, such narratives have to be left out of consideration from the outset as possess no feature that makes them intelligible to the Primary Grade. The test of their acceptibility is not that no unassimilated fragment remains; for in that case few would be left for even the Third Course; but those should be excluded for which the apperceptive fulcrum is altogether lacking in the concept world of little children (cf. p. 253). Accordingly all stories of the Old Testament should be left out of consideration which deal with the people of Israel in general; for the term “people” is foreign to the concept world of little children. But the family stories and episodes from the lives of eminent men should occupy the foreground; for it is in families that children originate, and even for the comprehension of a biography the requisite apperceptive fulcrum is found in the child’s concept sphere (cf. Schueren, op. cit., p. 10). Accordingly certain sections of the New Testament, such as Zacharias’ hymn of thanksgiving, Jesus’ talk with Nicodemus, the conversation of Jesus with the Samaritan woman, Jesus’ discourses on the end of the world, etc., are excluded from the outset as material for the Second Course. We have come to the limit of the efficiency of this agency, and it is well that we should recognize that fact.

Out of the question for us is the haphazard selection
\quellseite{445}
of the lessons now from this period, now from that, after the manner of the old lesson plans of the I.S.S.A. The stories of the Bible should be presented in large groups, in coherent, integral masses of thought, derived from no more than two, or at the highest, three periods. Only thus can the children really penetrate into the material; only thus is it possible for those Bible characters to become factors in the child’s inner life and stir his varied feelings. Splendid is what Fankhauser says upon this subject: “When we have begun to tell of some personage, for instance, Abraham, the child has taken an interest in him if the proper method has been pursued. There is a man who has gone from home, away from father, mother, friends, and relatives, without knowing what is to come of it, blindly following the command of God. Something has been wrong in our presentation if the children do not desire to find out from us the next time how that man who has won their sympathy, fared further. It is a fine thing, then, to be able to meet the desire of the child; that means we have gained a great advantage for our instruction at the very outset. We do not have to explain at length whom we are going to talk about, in what circumstances he lived, etc. — things surely necessary where an unknown person is to be introduced. It is just like a piece of embroidery work: mother has already placed it on hoops and threaded the needle; a few feathers of a many-colored tropical bird begin to appear in the fabric; and wondering what the figure will look like when it is completed, we anxiously watch mother’s moves till all the gorgeous colors are blended in the finished product” (\emph{Die biblische Geschichte in Sonntagsschule und Religionsstunde}, p. 18). Another thing, but one new story is to be selected for each Sunday. It is true that there are still Biblical Histories which harness two or three stories mechanically together in order to produce 2x52 lessons, according to the example of old
\quellseite{446}
Huebner. From the pedagogic standpoint, this is illsuited even to the parish school, much more so to the Sunday School where the “driver’s stick” must be still less in evidence. “Not too much should be given at one time; but what is given should not be a fragment but a coherent, well-rounded unit” (Schueren, \emph{op. cit.}, p. 17). How could anyone possibly think of launching forth a second story immediately after the first while the first story still holds the mind of the child enthralled and is only just about to find its way to his heart, especially if the second one portrays a quite different picture and has a different aim? The effect of the second picture would merely be to obliterate or efface the first; it would nullify the intended effect upon the mind (cf. p. 207). That proceeding would tend to make the children superficial of mind and dull of soul. Is it not the imperative duty of the teacher to focus the attention of his pupils on one object until a clear intuition or a clear concept has been formed? And does he not conscientiously try to fulfill this duty even if the lightning like rapidity with which his pupils’ attention shifts from one object to another, renders it now and then very difficult for him to do it? Does he not at all times keep in mind the law of the narrowness of consciousness? How, then, should he suddenly become a merciless driver, chasing the poor children from one story to another and never permitting the concepts that have been aroused, to take firm root in the soul! (Cf. Fankhauser, \emph{op. cit.}, p. 21.) Only if the several stories constitute a catechetical unit, may they be presented in one lesson.

But what order is to be observed in the selection of the several stories? He who knows what a precious aid the church year is in the inculcation of the essential facts of salvation (cf. p. 375), and for this reason is concerned about introducing the young to the idea of the church year and preparing them, also in this direction, for intelligent participation
\quellseite{447}
in the worship of the mature congregation, will follow the order of the church year also in the Sunday School, and always present, at least during the festive half of the church year, stories from the New Testament, with special emphasis upon the outstanding events in the life of Jesus; while the other half of the church year may be devoted to the story of the Patriarchs, of Moses, etc., that is such units as are required by the curriculum based upon the successive cultural epochs. — Finally, because the Sunday School, especially where there is no parish school, has to be a school rather than a service of worship, and because we do not merely wish to arouse a transient emotion in the soul, but to impart a treasure for life, care must be taken that familiarity with the subject-matter is attained. To this end there will not only be a review every seventh or eighth Sunday, but there will be also a resumption of the same topics every third or fourth year. This, from another point of view, brings us back to the requirement: the three or two grades of the Second Course have the same lesson story though adapted to the various stages of development represented by the several grades. If the objection is raised that the recurrence of the same material every third year kills interest, we reply with Schueren: “He who thinks that the repeated relation of a story causes ennui, is not acquainted either with the character of the Bible story or of the child. Where the real story-telling note is struck even in a moderate measure, the biblical story will always appear new to the children; with increasing interest they will listen to the person telling a story that has caused them joy once before. How is a mother who has the story-telling gift to any degree whatever, tormented by the children to tell the same story for the tenth or twentieth time! How the children hang upon her lips when she tells the story! It impresses them as if it had never been told! Teachers whose experience is different should not accuse the children but themselves” (\emph{op. cit.}, p. 10).

\quellseite{448}
In the \emph{Wartburg Lesson Helps}, the following stories constitute the Second Course: First Year: 1. How God created the world.—2. How God made man happy.—3. How man sinned.—4. Through sin man was lost and condemned.—5. How the Savior came into the world.—6. The child Jesus is brought to the temple.—7. Wise men from the East come to the child.—8. How the child Jesus was saved from great danger.—9. Jesus, when a boy twelve years of age, again comes to the temple.—10. How Jesus changed water into wine at the wedding at Cana.—11. How Jesus blessed Peter’s draught of fishes.—12. How Jesus saved His disciples during the storm.—13. Jesus feeds the hungry multitude.—14. How Jesus healed the man sick of palsy both in body and in soul.—15. How Jesus comforted a sorrowing father by restoring his daughter to life.—16. Jesus desires that little children should be brought to Him.—17. Jesus is our Good Shepherd.—18. How Jesus died for us on Good Friday.—19. How Jesus arose from the dead and appeared to Mary Magdalene.—20. Jesus would help the greatest of sinners, if they would only come to Him.—21. Jesus tells who the neighbor is, whom we should help.—22. Jesus tells us to forgive those who sin against us.—23. Jesus tells us to pray.—24. Jesus ascends into heaven in order to be with us always.—25. Jesus sends the Holy Spirit, who shall lead us to Jesus.—26. Of godly Abraham who by faith left his home.—27. Abraham would rather separate from Lot than live in strife with him.—28. How Abraham proved to be a friend in need.—29. What pleased God most in Abraham.—30. How Abraham spake with God as a friend speaks with a friend.—31. How Sodom and Gomorrah were destroyed because of their sin.—32. God gives Abraham a son, but Abraham loves God more than his son.—33. How God provided a wife for Isaac.—34. How Jacob deceived his father and his brother.—35. How Jacob fled from his brother and in his dream saw the heaven open.—36. How God prospered Jacob in the strange land and safely brought him home again.—37. Of Joseph, Jacob’s dearest son.—38. How Joseph was sold by his brethren for twenty pieces of silver.—39. How Joseph was humbled and imprisoned.—40. How Joseph was exalted and made lord of all Egypt.—41. Joseph’s brethren come to Egypt for the first time.—42. Joseph’s brethren come to Egypt for the second time.—43. Joseph makes himself known to his brethren.—44. Joseph cares for his father in Egypt.

\emph{The Second Year:} 1. How sin came into the world.—2. How sin spread so rapidly.—3. How God destroyed man because of sin.—4. How the Saviour’s messenger was born.—5. How the Saviour came into the world.—6. How John prepared the way for Jesus and baptized Him.—7. How Jesus resisted the devil’s attempt
\quellseite{449}
to dissuade Him from being our Saviour.—8. How Jesus healed the centurion’s servant, and encountered great faith.—9. How Jesus restored a son to his mother.—10. How Jesus heard the prayer of a heathen woman.—11. How Jesus entered the house of a sinner and made a new man out of him.—12. How Jesus entered Jerusalem on Palm Sunday.—13. How the Lord’s Supper was instituted.—14. How hard it was for Jesus to suffer and die for us.—15. How Jesus was betrayed by one of His disciples.—16. How one of His disciples denied Jesus.—17. How Jesus was carried from one court to another on account of our sin.—18. How Jesus died and was buried on Good Friday.—19. How Jesus arose from the dead and appeared to Mary Magdalene.—20. How Jesus comforted two disciples on their ways to Emmaus.—21. How Jesus appeared to His disciples and brought Thomas to believe in His resurrection.—22. How Jesus prepared salvation for all, though believers only partake in it.—23. How he only will be saved who perseveres in faith.—24. How Jesus ascends into heaven to be with us always.—25. How Jesus sends the Holy Spirit that He may lead us to Christ.—26. How a small boy was wonderfully preserved and received the name Moses.—27. How God appeared to Moses and called him to be the saviour of the people of Israel.—28. What Moses accomplished before proud Pharaoh.—29. How Moses led the people of Israel through the Red Sea.—30. How God cared for His people on their journey.—31. How God gave the Ten Commandments to the people of Israel.—32. How Moses took leave of the people of Israel and died.—33. How Joshua brought the people of Israel into the land of Canaan.—34. Of Eli’s wicked sons and pious Samuel.—35. How Samuel appointed Saul to become king.—36. How God rejected Saul because of his disobedience.—37. How God selected pious David to be king.—38. How David trusted in God and slew Goliath.—39. How David suffered much from Saul, though he found a good friend in Jonathan.—40. How David refrained from rendering evil for evil.—41. How God exalted David and made him king.—42. How Absalom sorely grieved David.—43. How Solomon, a king after God’s heart, built the temple in Jerusalem.—44. How God permitted a famine to come over Israel and miraculously fed the prophet Elijah.

The First Course or Beginners’ Department comprises the children from five to seven years of age. No regular lesson can be assigned them. They must be entertained rather than instructed, yet so entertained that they are led into an understanding of religious fundamentals, and the new life implanted
\quellseite{450}
in Baptism is nourished and made to become active. In the author’s book, For Beginners, it is practically worked out, how these children are to be instructed. It contains the following lessons:

1. God made all things.—2. God made Adam, the first man.—3. What God has given me.—4. God cares for me.—5. God protects us by His angels.—6. God protects us by His angels.—7. How God has made me His child.—8. How we may speak with our Heavenly Father.—9. How we may speak with our Heavenly Father.—10. God hears our prayers.—11. Jesus is our Good Shepherd.—12. God sees and knows everything.—13. God knows and sees everything.—14. How God punishes us if we do evil.—15. How God punishes us if we do evil.—16. God keeps His promise.—17. We should thank God when He has helped us.—18. God looks upon our hearts.—19. What we should do on Sunday.—20. We should help all those who are in need.—21. The Bible is the best book.—22. How beautiful it is in heaven.—1. Christmas is coming.—2. An angel comes to the virgin Mary.—3. God gives the child Jesus to the virgin Mary.—4. The angels come to the shepherds.—5. The shepherds come to see the child Jesus.—6. How the shepherds made known the birth of Jesus.—7. The wise men seeking Jesus.—8. The wise men find and worship Jesus.—9. God protects the child Jesus from Herod.—10. How God and man are pleased with the child Jesus.—11. The boy Jesus goes to celebrate Easter.—12. The boy Jesus in the temple hears the word of His father in heaven.—13. Mary finds the boy in the temple.—14. Jesus finds His first disciples.—15. How Jesus changes water into wine.—16. Jesus feeds the hungry.—17. The winds and the seas obey Jesus.—18. Jesus heals the sick man.—19. Jesus opens the eyes of two blind men.—20. Jesus raises the daughter of Jairus to life.—21. How the Lord Jesus blessed children.—22. Jesus enters Jerusalem on Palm Sunday.—23. How Jesus prayed in the Garden of Gethsemane.—24. Jesus is taken prisoner.—25. The Lord Jesus is crowned with thorns.—26. Jesus bears His cross to Golgotha.—27. Jesus dies on the cross for us.—28. The Lord Jesus is buried.—29. Jesus is raised from the dead and shows Himself to Mary Magdalene.—30. Jesus is taken up into heaven.

The Third Course, or Senior Department (early and later adolescence) of the Sunday School is chiefly concerned with introducing the pupils to the Bible, to explain it, and to train for a personal use of it. Where introduction to the history of the Church and her labors is not assigned to the Young
\quellseite{451}
People’s Society, the place for that, too, would be here. At the very outset it should be borne in mind that a dry introduction to the extent of the canon, the series of books composing it, details as to the origin of the canon, and the history of the Bible—all necessary things to be known—accomplishes very little after all. The young must be led into the Scriptures themselves, they should read them again and again. While the necessary explanations should be forthcoming in connection with the reading of the Scriptures, they should not bear the character of an exegesis, such as would be in place in a theological seminary; but it should most carefully bring out the essentials and fundamentals, show their bearing upon the lessons of the Catechism, and always aim at practical results in the lives of the young people. Whether it may not be expedient here and there to take a shorter route through the whole Bible at the outset, to be followed by a more thorough course afterward, depends upon circumstances.

The Saturday School should go hand in hand with the Sunday School, paying particular attention to the “penetration” and “application” of the Bible stories which have been “presented” in the Sunday School. How that is to be done, can be seen from § 34. Here time is afforded for an analysis of the Catechism, for the necessary explanation of its terms, and a preparation by these means for more thorough subsequent study of the Catechism, always, be it remembered, in connection with Biblical History. Here is also the occasion for practicing church hymns and for their preliminary explanation. We know from practical experience how, in this way, the pastor gets the Sunday School really under his control and effectually supplements its work.

The great value of the Summer School is found in this that it affords a daily opportunity for handling religious material with the pupils. The teacher is enabled to introduce his pupils to larger complexes; to delineate with them characters,
\quellseite{452}
that of Abraham, Jacob, Joseph, Moses, David; to survey the unfolding of the prophetic ministry of Christ: only thus is it possible for the child to identify himself with another age and another character; only thus it is possible for him to enter upon an “ideal association” with the persons of sacred history. Separate instruction in the Catechism should be given in this school to the older pupils. The Ten Commandments especially need to be stressed in this connection, in order to acquaint the children with the order which the Father in heaven wants observed in His house (cf. p. 313). If it be true that the explanations of the Catechism should never smack of doctrinal barrenness, but always throb with life, this is true of the Commandments in particular. Instead of insisting upon definitions of the many conceptions occurring in the First Chief Part, the children should be made to visualize pictures from life; there they will be shown in what Christian life and conduct consists. Alongside the First Chief Part, the Second might receive attention, but only in outline; it is sufficient so to explain it that it can be intelligently memorized. Narratives from the history of the Church and Missions, the fundamentals of Biblical Geography, the practise and discussion of the more important church hymns—these are subjects that will be used to supplement the lesson material so far mentioned.

The instruction of catechumens is bound to have its center in the Catechism. If a parish school is available and thorough instruction has been given in the first three Chief Parts, it is practicable in these circumstances to abandon the regular order of the Five Chief Parts and to select another way. Confirmation, for which the catechumens are to be prepared, may be made the starting point; from there one may go back to Baptism, proceed to the Creed next, and finally, at the hand of the Decalogue and the Lord’s Prayer, describe the Christian’s walk, the whole process to issue in Holy Communion and Confession. Where there is a thorough knowledge of Biblical History as
\quellseite{453}
well, the following process, too, is feasible: Creation and the fall of man may be made the starting point, with which the First Article readily connects itself, whereupon, after a brief glance at the redemptive purpose (cf. Gen. 3:15), the story of salvation down to the resurrection of Christ may be gone through with the pupils (cf. pp. 299–307). Here a stop is made and the Second Article taken up and the story of salvation pursued down to Pentecost, whereupon the Third Article is explained and the Sacrament of Baptism. Then, at the hand of the Decalogue and the Lord’s Prayer, the conduct of the Christian may be described, with special reference to the example of Christ, and instruction would conclude with an explanation of Confession, the Holy Supper, and Confirmation. Where neither of these postulates exist, the instruction of catechumens will closely follow the parts of the Catechism: it will bring out particularly the essentials, with a constant view to the practical Christian life and the spirit permeating it, a pastoral atmosphere pervading the whole process. Such instruction will further make a beginning of introducing the pupil to the Scriptures, cast light upon Lutheran usage, the church year, the order of service, and give practical information regarding the tasks of the Christian congregation (especially missions). Compare O. Hardeland, \emph{52 Konfirmandenstunden}, ⁴1910; R. Steinmetz, \emph{Das gute Bekenntnis, Hilfsbuch zur Vorbereitung auf die Konfirmation}; Steinmetz, \emph{Die Bereitung zur Konfirmation in Lehre und Leitung}, 1910; Vogel, \emph{Seelsorgerlicher Konfirmandenunterricht}, 1911.

In the Bible Class of the Sunday School (cf. what was said above) and in the Young People’s Society, the object must be an introduction to the Scriptures and to the secondary material of instruction as it was called in § 29. How to introduce the pupils to the Holy Scriptures, has already been indicated (p. 372); it is practically demonstrated in the Third Course of the Wartburg Lesson Helps (Book of Books; Book
\quellseite{454}
of Life), composed by the author of this textbook. How the secondary material is to be used in Young People’s Societies we show by adding several lesson plans. In regard to presentation compare Part V. In part also this material has been worked out practically, for instance in the author’s book, \emph{\emph{The Life of Dr. Martin Luther, Sketched for Young People’s Societies, and the Necessary Directions for General Discussion Appended}}, 1917, and in the department \emph{Luther (or Wartburg) League Topics} in the Lutheran Herald.

1. The Church and Her History: The Heathen World and its Moral Condition at the Time of Christ; The Preaching of the Gospel Among the Heathen Through the Apostles; The Persecution of Christians under Nero; Polycarp and Justin; The Persecution in Gaul; Perpetua and Felicitas; Persecution under Diocletian; Constantine and the Victory of Christianity; Julian the Apostate; Pictures from the Life of the Early Church; Athanasius and the Doctrinal Conflicts of his Time; Antonius and Monasticism; Chrysostom and his Congregation in Constantinople; The Origin of Papacy; Mohammed and the Judgment of God upon the Church; The Gospel in Germany; Charlemagne; Missions North and East; Gregory VII and Henry IV; The Crusades; The Monastic Orders; The Decay of the Church; The Waldensians; Wiclif and Hus.—Luther; Melanchthon; Zwingli and Calvin; The Gospel in England; The Gospel in France; The Counter-Reformation and the Thirty Years’ War; Paul Gerhardt; Duke Ernest of Gotha; Phil. Spener; Rationalism; The New Spring; Origin of the Union.—Muehlenberg and the Founding of the Lutheran Church in America; Rationalism in the Lutheran Church of America and the Founding of the General Synod; Founding of the Missouri Synod; Founding of the Iowa Synod; Founding of the General Council; Task of the Lutheran Church of America.

—2. The Church and Her Confessions: The Duty of the Church to Confess her Faith and the Three Oecumenical Confessions; The Lutheran Church and her Confessions; The Confessions of the Lutheran Church and the Reformed Church; The Lutheran Confession and the Union (with special reference to its representatives in America); The Lutheran Confession and Methodism; The Lutheran Confession and the Baptists; The Lutheran Confession and the Adventists; The Lutheran Confession and Christian Science.—3. The Church and her Problems: Church and World; Church and Politics; The Church and Political
\quellseite{455}
Parties; The Church and the Secular Press; The Church and the Labor Movement; The Church and Capital; The Church and Prohibition; The Church and the Lodge; The Church and the Public School; The Church and the Language Question.—4. The Church and Her Tasks: a) Foreign Missions: The Duty of Carrying on Foreign Missions; Ziegenbalg and Schwarz, and Lutheran Missions in India; Hans Egede the Apostle of Greenland; Zinzendorf and the Moravian Missions; Williams the South Sea Apostle; W. Carey the Missionary to the Bengalese; Guetzlaff the Founder of the Chinese Missions; The Leipzig Missions in India; L. Harms and the Hermannsburg Missions in Africa; Livingstone the Explorer and Missionary in the Dark Continent; Madagascar, the Grief and the Joy of Missions; The New Dettelsau Missions in Australia and New Guinea; Baierlein and Lutheran Missions Among the Indians; A Present Survey of Foreign Missions.—b) Inner Missions: Inner Missions the Duty of the Church of Jesus; A. H. Francke and Inner Missions; Muehlenberg’s Labors Among the American Dispersion; Wyneken’s Pioneer Labors; W. Loehe and His Efforts on behalf of the Lutherans in America; The Pioneers of the Iowa Synod; The Organization of a Mission Congregation.—W. Loehe and the Deaconess Work; Theodore Fliedner and Kaiserswert; Wichern and the Rauhe Haus in Hamburg; G. Mueller the Orphan Father of Bristol; Survey of Inner Missions in Germany; Passavant the Father of Lutheran Inner Missions in America; A Day in the Deaconess Home in Philadelphia or Milwaukee; A Day in the Orphans Home of Mount Vernon; A Day in the Orphans Home of Waverly; City Mission in Minneapolis; Our Homes for the Aged; Our Higher Institutions.—5. The Church and Her Usages: The Necessity for Fixed Usages in the Church; The Church Year; The Pericopes; The Celebration of the Sunday; The Order of the Regular Sunday Service; The Order of the Communion Service; The Order of the Baptismal Service; Announcement for Confession; Sponsorship; System in Giving.—6. The Church and Her Constitution: The Proper Organization of the Local Lutheran Congregation; The Ideal Constitution for the Local Congregation; The Necessity for Membership in a Local Congregation; Church Discipline; The Synod and Its Organization; Connection Between the Synod and the Congregation; The Correct Relation to other Synods and Its Congregations and Members.

In addition to the several series given above, such as these would be feasible: Bible Characters: Adam and Eve; The Sons of Adam
\quellseite{456}
and Eve, or The First Stage of Culture; Noah and the Flood; Nimrod and the Origin of World Power and of Paganism; Abraham; Isaac; Jacob; Joseph; Moses; Joshua; Samuel; David; Solomon; Elijah; Isaiah; Jeremiah; Daniel; Ezra; The Maccabees; Herod.—Jesus’ Life in the Light of the Ten Commandments; Mary the Mother of the Lord; Mary the Sister of Lazarus; Mary Magdalene; Peter; John; Judas; Stephen; James the Just; The First Christian Congregation; Ananias and Sapphira; Paul on the Way to Damascus; Paul in Lystra; Paul in Athens and Corinth; Paul in Caesarea; Paul in Rome; Timothy (here the sections of Scripture bearing upon the subject should be read by the youth by way of preparation).—Church Hymns Explained from the Aspect of Their Origin: The Church Hymn and Its Significance in Reformation Times: “A Mighty Fortress” and Parallels from Luther’s Life; Speratus and “To Us Salvation Now Has Come”; Ph. Nicolai and “Wake, Arise, the Voice Is Calling”; M. Rinkart and “Now Thank We All Our God”; P. Gerhardt and “Thy Way and All Thy Sorrows”; Neumark and “If Thou but Suffer God to Guide Thee”; Fleming and “In All that may Betide me, The Most High God Shall Guide Me” (In allen meinen Taten); John Heermann the Singer in the School of Suffering; John Frank the Godly Burgomaster and his Hymn “Deck Thyself, O Soul, with Gladness”; Gellert and “Jesus Lives! No Longer Now”; J. Sebastian Bach the Godly Organist; “Jesus Christ, My Sure Defense” and the Blessings in Its Wake.—Biblical, Ethical, and Apologetic Materials: The Old Testament and Its Origin; The New Testament and Its Origin; The History of the Bible in the Church; The Psalms and Their Significance for the Inner Life; The Prophetic Writings and their Significance for the Inner Life; How the Bible Should Be Read to Build up the Inner Life.—The Blessing of True Friendship; Youth’s Happy Time; The Special Dangers of Youth; The Necessity of Conversion; Prayer and Its Blessing; The Highest Blessings of All; Conscience and Its Cultivation; Truthfulness a Fundamental Condition of True Piety; Healthy Recreation; Getting Ready for the Last Journey; Our Duty Toward the Dead; the Three Articles of Our Christian Faith and the Intellectual Movements of the Present; Is It True that Only Fools Believe?—Pictures from the Holy Land might similarly do service.
